\documentclass[../../main.tex]{subfiles}

\begin{document}

\section{线性泛函、有界线性泛函以及弱拓扑}\label{section:ch14-线性泛函弱拓扑}

令 $X$ 为线性空间. 我们把 $X$ 上的线性实值函数的线性空间记为 $X^{\#}$. 对 $\psi\in X^{\#}$, $\psi\neq 0$ 及 $x_{0}\in X$, 使得 $\psi(x_{0})\neq 0$, 我们宣称 $X$ 可被表示为直和:
\begin{align}
X=[\ker\psi]\oplus\operatorname{span}[x_{0}]\label{align:ch14-直和核}
\end{align}
其中 $\psi$ 的核 $\ker\psi$ 是子空间 $\{x\in X\mid\psi(x)=0\}$. 事实上, 显然, $\ker\psi\cap\operatorname{span}[x_{0}]=\{0\}$. 另一方面, 我们可以将每个 $x\in X$ 写为
\begin{align*}
x=\left[x-\frac{\psi(x)}{\psi(x_{0})}\cdot x_{0}\right]+\frac{\psi(x)}{\psi(x_{0})}\cdot x_{0},\quad\text{且有 }\psi\left(x-\frac{\psi(x)}{\psi(x_{0})}\cdot x_{0}\right)=0.
\end{align*}
观察到对实数 $c$, 若 $x_{0}$ 属于 $X$ 且 $\psi(x_{0})=c$, 则
\begin{align*}
\psi^{-1}(c)=\{x\in X\mid\psi(x)=c\}=\ker\psi+x_{0}.
\end{align*}
因此, 由\eqref{align:ch14-直和核}, 若 $X$ 是维数为 $n$ 的有限维空间而 $\psi$ 非零, 则对每个 $c\in\mathbf{R}$, 水平集 $\psi^{-1}(c)$ 是 $X$ 的 $(n-1)$ 维子空间的平移.

若 $X$ 的线性子空间 $X_{0}$ 具有以下性质: 存在某个 $x_{0}\in X$, $x_{0}\neq 0$, 使得 $X=X_{0}\oplus\operatorname{span}[x_{0}]$, 则 $X_{0}$ 称为在 $X$ 中的余维是 $1$. 余维是 $1$ 的子空间的平移称为超平面.

\begin{proposition}\label{proposition:ch14-余维1}
线性空间 $X$ 的线性子空间 $X_{0}$ 余维是 $1$ 当且仅当对某个非零 $\psi\in X^{\#}$, $X_{0}=\ker\psi$.
\end{proposition}

\begin{proof}
我们已观察到非零线性泛函的核的余维是 $1$. 反过来, 假定 $X_{0}$ 是余维为 $1$ 的子空间. 则存在向量 $x_{0}\neq 0$ 使得 $X=X_{0}\oplus\operatorname{span}[x_{0}]$. 对 $\lambda\in\mathbf{R}$ 与 $x\in X_{0}$, 定义 $\psi(x+\lambda x_{0})=\lambda$. 则 $\psi\neq 0$, $\psi$ 是线性的且 $\ker\psi=X_{0}$.
\end{proof}

以下命题告诉我们线性空间上的线性泛函是充足的.

\begin{proposition}\label{proposition:ch14-延拓线性泛函}
令 $Y$ 为线性空间 $X$ 的线性子空间. 则 $Y$ 上的每个线性泛函可延拓为整个 $X$ 上的线性泛函. 特别地, 对每个 $x\in X$, $x\neq 0$, 存在 $\psi\in X^{\#}$ 使得 $\psi(x)\neq 0$.
\end{proposition}

\begin{proof}
如同我们在前一章观察到的(见该章的习题 36), $Y$ 在 $X$ 中有线性补, 即存在 $X$ 的线性子空间 $X_{0}$, 使得以下直和分解
\begin{align*}
X=Y\oplus X_{0}
\end{align*}
成立. 令 $\eta$ 属于 $Y^{\#}$. 对 $x\in X$, 我们有 $x=y+x_{0}$, 其中 $y\in Y$ 而 $x_{0}\in X_{0}$. 定义 $\eta(x)=\eta(y)$. 这定义了 $\eta$ 到整个 $X$ 上的线性泛函的延拓.

现在令 $x\neq 0$ 属于 $X$. 定义 $\eta:\operatorname{span}[x]\to\mathbf{R}$ 为 $\eta(\lambda x)=\lambda\cdot\|x\|$, 根据证明的第一部分, 线性泛函 $\eta$ 可延拓为整个 $X$ 上的线性泛函.
\end{proof}

我们对赋范线性空间 $X$ 以及 $X^{\#}$ 的包含于 $X$ 的对偶空间 $X^{*}$ 的子空间特别感兴趣, $X^{*}$ 即关于由范数诱导的拓扑连续的线性泛函的线性空间. 若 $X$ 是有限维赋范线性空间, 则 $X$ 上的每个连续线性泛函都属于 $X^{*}$(见习题 4). 该性质刻画了有限维赋范线性空间.

线性空间 $X$ 的子集 $\mathcal{B}$ 称为 $X$ 的 Hamel 基, 若 $X$ 中的每个向量可唯一地表示为 $\mathcal{B}$ 中向量的有限线性组合. 我们把用 Zorn 引理推出每个线性空间具有 Hamel 基留作练习(见习题 17).

\begin{proposition}\label{proposition:ch14-有限维X星}
令 $X$ 为赋范线性空间. 则 $X$ 是有限维的当且仅当 $X^{\#}=X^{*}$.
\end{proposition}

\begin{proof}
我们把以下证明留作练习: 由于有限维线性空间上的所有范数是等价的, 所以这样的空间上的所有线性泛函是有界的. 假设 $X$ 是无穷维的. 令 $\mathcal{B}$ 为 $X$ 的 Hamel 基. 不失一般性, 我们假设 $\mathcal{B}$ 中的向量是单位向量. 由于 $X$ 是无穷维的, 我们可以选取 $\mathcal{B}$ 的一个可数无穷子集, 我们把它列举为 $\{x_{k}\}_{k=1}^{\infty}$. 对每个自然数 $k$ 与向量 $x\in X$, 定义 $\psi_{k}(x)$ 为 $x$ 在 Hamel 基 $\mathcal{B}$ 下的展开式中 $x_{k}$ 的系数. 则每个 $\psi_{k}$ 属于 $X^{\#}$. 因此对所有 $x\in X$, 定义为
\begin{align*}
\psi(x)=\sum_{k=1}^{\infty}k\cdot\psi_{k}(x)
\end{align*}
的泛函 $\psi:X\to\mathbf{R}$ 也属于 $X^{\#}$. 该线性泛函不是有界的, 由于每个 $x_{k}$ 是单位向量且 $\psi(x_{k})=k\cdot\psi_{k}(x_{k})=k$.
\end{proof}

以下线性泛函的代数性质在建立弱拓扑的性质时是有用的.

\begin{proposition}\label{proposition:ch14-核交}
令 $X$ 为线性空间, 泛函 $\psi$ 属于 $X^{\#}$, 而 $\{\psi_{i}\}_{i=1}^{n}$ 包含于 $X^{\#}$. 则 $\psi$ 是 $\{\psi_{i}\}_{i=1}^{n}$ 的线性组合当且仅当
\begin{align}
\bigcap_{i=1}^{n}\ker\psi_{i}\subseteq\ker\psi.\label{align:ch14-核包含}
\end{align}
\end{proposition}

\begin{proof}
显然, 若 $\psi$ 是 $\{\psi_{i}\}_{i=1}^{n}$ 的线性组合, 则包含关系\eqref{align:ch14-核包含}成立. 我们归纳地证明反命题. 对 $n=1$, 假定\eqref{align:ch14-核包含}成立. 我们假设 $\psi\neq 0$, 否则没什么可证的. 选取 $x_{0}\neq 0$ 使得 $\psi(x_{0})=1$. 则 $\psi_{1}(x_{0})\neq 0$, 由于 $\ker\psi_{1}\subseteq\ker\psi$. 然而, $X=\ker\psi_{1}\oplus\operatorname{span}[x_{0}]$. 因此, 若我们定义 $\lambda_{1}=1/\psi_{1}(x_{0})$, 通过直接代换, 我们看到, $\psi=\lambda_{1}\psi_{1}$. 现在假设: 对任何线性空间上的线性泛函, 若\eqref{align:ch14-核包含}对 $n=k-1$ 成立, 则 $\psi$ 是 $\psi_{1},\cdots,\psi_{k-1}$ 的线性组合. 假定\eqref{align:ch14-核包含}对 $n=k$ 成立. 若 $\psi_{k}=0$, 没什么可证的. 因此选取 $x_{0}\in X$ 满足 $\psi_{k}(x_{0})=1$. 则 $X=Y\oplus\operatorname{span}[x_{0}]$, 其中 $Y=\ker\psi_{k}$, 因此,
\begin{align*}
\bigcap_{i=1}^{k-1}[\ker\psi_{i}\cap Y]\subseteq\ker\psi\cap Y.
\end{align*}
根据归纳假设, 存在实数 $\lambda_{1},\cdots,\lambda_{k-1}$ 使得
\begin{align*}
\text{在 }Y\text{ 上},\quad\psi=\sum_{i=1}^{k-1}\lambda_{i}\cdot\psi_{i}.
\end{align*}
直接代换表明, 若我们定义 $\lambda_{k}=\psi(x_{0})-\sum_{i=1}^{k-1}\lambda_{i}\cdot\psi_{i}(x_{0})$, 则
\begin{align*}
\text{在 }X\text{ 上},\quad\psi=\sum_{i=1}^{k}\lambda_{i}\cdot\psi_{i}.
\end{align*}
\end{proof}

回忆对集合 $X$ 上的两个拓扑 $\mathcal{T}_{1}$ 与 $\mathcal{T}_{2}$, 若 $\mathcal{T}_{1}\subseteq\mathcal{T}_{2}$, 我们说 $\mathcal{T}_{1}$ 比 $\mathcal{T}_{2}$ 弱或 $\mathcal{T}_{2}$ 比 $\mathcal{T}_{1}$ 强. 观察到 $X$ 上的函数关于 $X$ 上的某个拓扑连续, 则它关于 $X$ 上的任何更强的拓扑连续, 但可能不关于某个较弱的拓扑连续. 若 $\mathcal{F}$ 是集合 $X$ 上的任何实值函数族, $X$ 上的由 $\mathcal{F}$ 诱导的弱拓扑, 或者说 $X$ 上的 $\mathcal{F}$ 弱拓扑, 定义为使得 $\mathcal{F}$ 中的每个函数连续的最弱拓扑(具有最少个数的集合的拓扑). $X$ 上的 $\mathcal{F}$ 弱拓扑在 $x\in X$ 的基由形如
\begin{align}
\mathcal{N}_{\varepsilon,f_{1},\cdots,f_{n}}(x)=\{x'\in X\mid|f_{k}(x')-f_{k}(x)|<\varepsilon,1\leqslant k\leqslant n\}\label{align:ch14-F弱基}
\end{align}
的集合组成, 其中 $\varepsilon>0$ 而 $\{f_{k}\}_{k=1}^{n}$ 是 $\mathcal{F}$ 的有限子族. 对集合上的拓扑, 我们知道集合中的序列关于该拓扑收敛于该集合的某个点意味着什么. 容易看出 $X$ 中的序列 $\{x_{n}\}$ 关于 $\mathcal{F}$ 弱拓扑收敛到 $x\in X$ 当且仅当
\begin{align}
\text{对所有 }f\in\mathcal{F},\quad\lim_{n\to\infty}f(x_{n})=f(x).\label{align:ch14-F弱收敛}
\end{align}
$X$ 上关于 $\mathcal{F}$ 弱拓扑连续的函数称为 $\mathcal{F}$ 弱连续. 类似地, 我们有 $\mathcal{F}$ 弱开集、$\mathcal{F}$ 弱闭集以及 $\mathcal{F}$ 弱紧集.

对线性空间 $X$, 考虑 $X$ 上由 $X^{\#}$ 的线性子空间 $W$ 诱导的弱拓扑是自然且有用的.

\begin{proposition}\label{proposition:ch14-W弱连续}
令 $X$ 为线性空间而 $W$ 是 $X^{\#}$ 的子空间. 则线性泛函 $\psi:X\to\mathbf{R}$ 是 $W$ 弱连续的当且仅当它属于 $W$.
\end{proposition}

\begin{proof}
根据 $W$ 弱拓扑的定义, $W$ 中的每个线性泛函是 $W$ 弱连续的. 剩下来要证明反命题. 假定线性泛函 $\psi:X\to\mathbf{R}$ 是 $W$ 弱连续的. 根据 $\psi$ 在 $0$ 的连续性, 存在 $0$ 的邻域 $\mathcal{N}$, 使得若 $x\in\mathcal{N}$, $|\psi(x)|=|\psi(x)-\psi(0)|<1$. $W$ 拓扑在 $0$ 的基有一个包含于 $\mathcal{N}$ 的邻域. 选取 $\varepsilon>0$ 与 $W$ 中的 $\psi_{1},\cdots,\psi_{n}$, 使得 $\mathcal{N}_{\varepsilon,\psi_{1},\cdots,\psi_{n}}\subseteq\mathcal{N}$. 因此,
\begin{align*}
\text{若对所有 }1\leqslant k\leqslant n,\quad|\psi_{k}(x)|<\varepsilon,\quad\text{则 }|\psi(x)|<1.
\end{align*}
根据 $\psi$ 与这些 $\psi_{k}$ 的线性, 我们有包含关系 $\bigcap_{k=1}^{n}\ker\psi_{k}\subseteq\ker\psi$. 根据\cref{proposition:ch14-核交}, $\psi$ 是 $\psi_{1},\cdots,\psi_{n}$ 的线性组合. 因此, 由于 $W$ 是线性空间, $\psi$ 属于 $W$.
\end{proof}

以上命题建立了 $X^{\#}$ 的线性子空间与由这样的子空间诱导的 $X$ 上的弱拓扑的一一对应.

\begin{definition}[弱拓扑]\label{definition:ch14-弱拓扑}
令 $X$ 为赋范线性空间. 由对偶空间 $X^{*}$ 诱导的 $X$ 上的弱拓扑称为 $X$ 上的弱拓扑.
\end{definition}

$X$ 上的弱拓扑在 $x\in X$ 的基由形如
\begin{align}
\mathcal{N}_{\varepsilon,\psi_{1},\cdots,\psi_{n}}(x)=\{x'\in X\mid|\psi_{k}(x'-x)|<\varepsilon,1\leqslant k\leqslant n\}\label{align:ch14-弱基}
\end{align}
的集合组成, 其中 $\varepsilon>0$ 而 $\{\psi_{k}\}_{k=1}^{n}$ 是 $X^{*}$ 的有限子族. 对与弱拓扑有关的拓扑概念, 我们用形容词"弱", 因此我们有弱紧集、弱开集等. 因此 $X$ 中的序列 $\{x_{n}\}$ 弱收敛到 $x$ 当且仅当
\begin{align}
\text{对所有 }\psi\in X^{*},\quad\lim_{n\to\infty}\psi(x_{n})=\psi(x).\label{align:ch14-弱收敛}
\end{align}
在 $X$ 中, 可以很方便地把 $X$ 中的序列 $\{x_{n}\}$ 弱收敛于点 $x\in X$ 写为 $\{x_{n}\}\to x$.

通常, 对赋范线性空间, 我们称由范数诱导的拓扑为 $X$ 上的强拓扑. 对 $X$ 是赋范线性空间而 $W$ 是 $X^{*}$ 的子空间, 在 $X$ 上的拓扑中存在着以下包含关系:
\begin{align*}
X\text{ 上的 }W\text{ 弱拓扑}\subseteq X\text{ 上的弱拓扑}\subseteq X\text{ 上的强拓扑}.
\end{align*}
我们从\cref{proposition:ch14-W弱连续}推出 $W$ 弱拓扑与弱拓扑重合当且仅当 $W=X^{*}$. 此外, 弱拓扑与强拓扑重合当且仅当 $X$ 是有限维的(见习题 6). 与赋范线性空间相应的拓扑概念若没有附加形容词, 则隐含假设所涉及的拓扑是强拓扑.

对于作为对偶空间的赋范线性空间, 存在除弱拓扑与强拓扑外的第三个重要的拓扑. 事实上, 对赋范线性空间 $X$ 与 $x\in X$, 我们定义泛函 $J(x):X^{*}\to\mathbf{R}$ 为
\begin{align*}
\text{对所有 }\psi\in X^{*},\quad J(x)[\psi]=\psi(x).
\end{align*}
显然, 估值泛函 $J(x)$ 是线性的且在 $X^{*}$ 上是有界的, 满足 $\|J(x)\|\leqslant\|x\|$. 此外, 算子 $J:X\to(X^{*})^{*}$ 是线性的, 因此 $J(X)$ 是 $(X^{*})^{*}$ 的线性子空间.

\begin{definition}[弱$^{*}$拓扑]\label{definition:ch14-弱星拓扑}
令 $X$ 为赋范线性空间. $X^{*}$ 上的由 $J(X)\subseteq(X^{*})^{*}$ 诱导的弱拓扑称为 $X^{*}$ 上的弱$^{*}$拓扑.
\end{definition}

$X^{*}$ 上的弱$^{*}$拓扑在 $\psi\in X^{*}$ 的基由形如
\begin{align}
\mathcal{N}_{\varepsilon,x_{1},\cdots,x_{n}}(\psi)=\{\psi'\in X^{*}\mid|(\psi'-\psi)(x_{k})|<\varepsilon,1\leqslant k\leqslant n\}\label{align:ch14-弱星基}
\end{align}
的集合构成, 其中 $\varepsilon>0$ 而 $\{x_{k}\}_{k=1}^{n}$ 是 $X$ 的有限子集. $X^{*}$ 的子集关于弱$^{*}$拓扑是开的, 称为是弱$^{*}$开的, 对其他拓扑概念类似处理. 因此 $X^{*}$ 中的序列 $\{\psi_{n}\}$ 弱$^{*}$收敛到 $\psi\in X^{*}$ 当且仅当
\begin{align}
\text{对所有 }x\in X,\quad\lim_{n\to\infty}\psi_{n}(x)=\psi(x).\label{align:ch14-弱星收敛}
\end{align}

简单来说, $X^{*}$ 中的弱$^{*}$收敛就是逐点收敛. 对赋范线性空间 $X$, $X^{*}$ 上的强、弱、弱$^{*}$拓扑通过以下包含关系相联系:
\begin{align*}
X^{*}\text{ 上的弱}^{*}\text{ 拓扑}\subseteq X^{*}\text{ 上的弱拓扑}\subseteq X^{*}\text{ 上的强拓扑}.
\end{align*}

\begin{definition}[自然嵌入]\label{definition:ch14-自然嵌入}
令 $X$ 为赋范线性空间. 线性算子 $J:X\to(X^{*})^{*}$ 定义为
\begin{align*}
\text{对所有 }x\in X,\ \psi\in X^{*},\quad J(x)[\psi]=\psi(x)
\end{align*}
称为 $X$ 到 $(X^{*})^{*}$ 的自然嵌入. 空间 $X$ 称为是自反的, 若 $J(X)=(X^{*})^{*}$. 习惯上记 $(X^{*})^{*}$ 为 $X^{**}$, 并称 $X^{**}$ 为 $X$ 的二次对偶.
\end{definition}

\begin{proposition}\label{proposition:ch14-自反弱星}
赋范线性空间 $X$ 是自反的当且仅当 $X^{*}$ 上的弱拓扑与弱$^{*}$拓扑是相同的.
\end{proposition}

\begin{proof}
显然若 $X$ 是自反的, 则 $X^{*}$ 上的弱与弱$^{*}$拓扑是相同的. 反过来, 假定这两个拓扑是相同的. 令 $\psi:X^{*}\to\mathbf{R}$ 为连续线性泛函. 根据弱拓扑的定义, $\psi$ 关于 $X^{*}$ 上的弱拓扑是连续的. 因此它关于弱$^{*}$拓扑是连续的. 我们从\cref{proposition:ch14-W弱连续}推出 $\psi$ 属于 $J(X)$. 因此 $J(X)=X^{**}$.
\end{proof}

目前, 我们未证实称 $J:X\to X^{**}$ 为"嵌入"的合法性, 因为我们还未证明 $J$ 是一对一的. 事实上, 我们甚至未证明在一般的赋范线性空间 $X$ 上存在任何非零的有界线性泛函. 对有界线性泛函我们需要命题 2 的一个变体, 即将介绍的 Hahn-Banach 定理提供这个变体且表明 $J$ 是等距. 当然, 我们已经研究了某些特殊的赋范线性空间的对偶空间. 例如, 若 $E$ 是实数的 Lebesgue 可测集而 $1\leqslant p<\infty$, Riesz 表示定理刻画了 $L^{p}(E)$ 的对偶.

\begin{exercise}\label{exercise:ch14-01}
(i)用 Zorn 引理证明每个线性空间有 Hamel 基.
(ii)证明无穷维 Banach 空间的任何 Hamel 基必须是不可数的.
(iii)令 $X$ 为定义在 $\mathbf{R}$ 上的所有多项式的线性空间. 证明不存在 $X$ 上的范数, 使得关于该范数 $X$ 是 Banach 空间.
\end{exercise}

\begin{exercise}\label{exercise:ch14-02}
在命题 4 的证明中验证两个关于直接代换的断言.
\end{exercise}

\begin{exercise}\label{exercise:ch14-03}
令 $X_{0}$ 为赋范线性空间 $X$ 的余维为 $1$ 的子空间. 证明 $X_{0}$ 关于强拓扑是闭的当且仅当对某个 $\psi\in X^{*}$, $X_{0}=\ker\psi$.
\end{exercise}

\begin{exercise}\label{exercise:ch14-04}
证明: 若 $X$ 是有限维赋范线性空间, 则 $X$ 上的每个线性泛函是连续的.
\end{exercise}

\begin{exercise}\label{exercise:ch14-05}
令 $X$ 为维数为 $n$ 的有限维赋范线性空间. 令 $\{e_{1},\cdots,e_{n}\}$ 为 $X$ 的一组基. 对 $1\leqslant i\leqslant n$, 定义 $\psi_{i}\in X^{*}$ 为: 对 $x=x_{1}e_{1}+\cdots+x_{n}e_{n}\in X$, $\psi_{i}(x)=x_{i}$. 证明 $\{\psi_{1},\cdots,\psi_{n}\}$ 是 $X^{*}$ 的一组基. 因此 $\dim X^{*}=n$.
\end{exercise}

\begin{exercise}\label{exercise:ch14-06}
令 $X$ 为有限维线性空间. 证明 $X$ 上的弱拓扑与强拓扑是相同的.
\end{exercise}

\begin{exercise}\label{exercise:ch14-07}
证明无穷维赋范线性空间的每个非空弱开子集关于该范数是无界的.
\end{exercise}

\begin{exercise}\label{exercise:ch14-08}
令 $X$ 为有限维线性空间. 证明自然嵌入 $J:X\to X^{**}$ 是一对一的. 接着用习题 5 证明 $J:X\to X^{**}$ 是映上的, 因此 $X$ 是自反的.
\end{exercise}

\begin{exercise}\label{exercise:ch14-09}
对于欧氏空间 $\mathbf{R}^{n}$ 的向量 $v\neq 0$, 直接给出线性泛函 $\psi:\mathbf{R}^{n}\to\mathbf{R}$ 使得 $\psi(v)=1$.
\end{exercise}

\begin{exercise}\label{exercise:ch14-10}
对于 $\ell^{2}$ 中的序列 $\{x_{n}\}\neq 0$, 直接给出连续线性泛函 $\psi:\ell^{2}\to\mathbf{R}$ 使得 $\psi(\{x_{n}\})=1$.
\end{exercise}

\begin{exercise}\label{exercise:ch14-11}
对 $L^{p}[a,b](1\leqslant p\leqslant\infty)$ 中的函数 $f\neq 0$, 直接给出连续线性泛函 $\psi:L^{p}[a,b]\to\mathbf{R}$ 使得 $\psi(f)=1$.
\end{exercise}

\begin{exercise}\label{exercise:ch14-12}
考虑具有最大值范数的 $C[a,b]$. 对于 $C[a,b]$ 中的函数 $f\neq 0$, 直接给出连续线性泛函 $\psi:C[a,b]\to\mathbf{R}$ 使得 $\psi(f)=1$.
\end{exercise}

\begin{exercise}\label{exercise:ch14-13}
对 $1<p<\infty$, 令 $Y$ 为 $L^{p}[a,b]$ 的余维为 $1$ 的闭子空间. 证明存在函数 $g\in L^{q}[a,b]$, 其中 $q$ 是 $p$ 的共轭, 使得
\begin{align*}
Y=\left\{f\in L^{p}[a,b]\mid\int_{[a,b]}f\cdot g\,dm=0\right\}.
\end{align*}
\end{exercise}

\begin{exercise}\label{exercise:ch14-14}
令 $X$ 为有限维赋范线性空间而 $\psi$ 属于 $X^{\#}\sim X^{*}$. 证明 $\ker\psi$ 关于 $X$ 中的强拓扑是稠密的.
\end{exercise}

\begin{exercise}\label{exercise:ch14-15}
令 $X$ 为限制在 $[a,b]$ 的多项式的赋范线性空间. 对 $p\in X$, 定义 $\psi(p)$ 为 $p$ 的系数的和. 证明 $\psi$ 是线性的. 若 $X$ 有由最大值范数诱导的拓扑, $\psi$ 是连续的吗?
\end{exercise}

\begin{exercise}\label{exercise:ch14-16}
令 $X$ 为仅有有限项非零的实数列的赋范线性空间. 对 $x=\{x_{n}\}\in X$, 定义 $\psi(x)=\sum_{n=1}^{\infty}x_{n}$. 证明 $\psi$ 是线性的. 若 $X$ 有由 $\ell^{\infty}$ 范数诱导的拓扑, $\psi$ 是连续的吗?
\end{exercise}

\begin{exercise}\label{exercise:ch14-17}
令 $X$ 为线性空间. $X$ 的子集 $E$ 称为是线性独立的, 若每个 $x\in E$ 不是 $E\sim\{x\}$ 中点的有限线性组合. 定义 $\mathcal{F}$ 为线性独立的 $X$ 的非空子集所组成的族. 用集合的包含关系将 $\mathcal{F}$ 排序. 用 Zorn 引理得出 $X$ 有 Hamel 基.
\end{exercise}

\begin{exercise}\label{exercise:ch14-18}
给出从赋范线性空间 $X$ 到赋范线性空间 $Y$ 的不连续线性算子 $T$ 的例子, 使得 $T$ 有闭的图像. (提示: 令 $\psi$ 为赋范线性空间 $X$ 上的不连续线性泛函, 而 $Y=\{y\in X\times\mathbf{R}\mid y=(x,\psi(x))\}$ 是 $\psi$ 的图像.)这与闭图像定理矛盾吗?
\end{exercise}

\end{document}